\documentclass[11pt]{amsart}

\usepackage[T1]{fontenc}
\usepackage{lmodern}
\usepackage{microtype}
\usepackage{mathtools,amssymb,amsthm}
\usepackage{array}
\usepackage[hidelinks]{hyperref}

\hypersetup{
  pdftitle={An Exhaustive Census Confirms a(6) = 1814 for Topologically Distinct
            Pizza Slicings}
}

\newtheorem{theorem}{Theorem}
\newtheorem{lemma}[theorem]{Lemma}
\newtheorem{corollary}[theorem]{Corollary}
\theoremstyle{definition}
\newtheorem{example}{Example}
\theoremstyle{remark}
\newtheorem{remark}[theorem]{Remark}

\newcommand{\Tn}{\mathcal T}
\newcommand{\Dsk}{\mathbb D}

\title[An exhaustive census confirms $a(6)=1814$]
{An Exhaustive Census Confirms $a(6) = 1814$ for Topologically Distinct
 Pizza Slicings}

\date{}

\subjclass[2020]{05A15, 52C30, 05C30}
\keywords{chord arrangement, pizza slicing, region graph, exhaustive enumeration,
line arrangement, stretchability, OEIS A272906}

\begin{document}

\begin{abstract}
Let $a(n)$ be the number of graphs, up to isomorphism, obtained by cutting a disk with
$n$ chords no three of which meet at a single strictly interior point, taking the pieces
as vertices and joining two pieces when a segment separates them; this is
OEIS \texttt{A272906}. That entry has carried the value $a(6)=1814$ since May 2016,
together with the note that the term is \emph{unconfirmed} because it was obtained from
guided random trials. We prove $a(6)=1814$. An exhaustive census of the
\emph{topological} six-chord arrangements in a disk produces exactly $1824$ types, whose
region graphs fall into exactly $1814$ isomorphism classes, giving $a(6)\le 1814$; and
$47$ explicitly printed twelve-tuples of integers, swept over all $10395$ perfect
matchings each, realize all $1824$ types by straight chords with rational endpoints,
giving $a(6)\ge 1814$. The value $1814$ is due to Jon Hart; we are unaware of a previous
exhaustive proof. A corollary of the realization step is that no non-stretchable
topological type occurs at $n=6$.
\end{abstract}

\maketitle

\section{The problem}

Let $\Dsk$ be a closed disk. A \emph{slicing} is a finite set of chords of $\Dsk$,
pairwise distinct, each with two distinct endpoints on the boundary circle, such that no
three of them pass through a common point of the open disk. The chords cut $\Dsk$ into
\emph{pieces}. The \emph{region graph} of the slicing has the pieces as vertices, two
pieces being adjacent when a line segment of positive length separates them; sharing a
single point, or sharing a boundary arc, does not count. Let
\[
 a(n)=\#\{\text{region graphs of slicings by $n$ chords}\}/{\cong}.
\]
This is OEIS \texttt{A272906} \cite{A272906}, \emph{Number of topologically-distinct
pizza slicings from $n$ chords in general position}, whose data line reads
\[
 1,\;1,\;2,\;5,\;19,\;130,\;1814
 \qquad (n=0,\ldots,6)
\]
and whose second comment records that the terms up to $a(5)=130$ are confirmed by
independent programs while $a(6)=1814$, obtained from guided random trials, is not.
Cutler, Karlsson and Sloane list the sequence in Open Problem~2 of \cite{CKS} and repeat
that the last term is unconfirmed.

\begin{theorem}\label{thm:main}
$a(6)=1814$.
\end{theorem}

\begin{corollary}\label{cor:stretch}
Every one of the $1824$ topological types of six-chord arrangements of a disk
$($Section~\textup{\ref{sec:upper}}$)$ is realized by six straight chords whose $12$
endpoints are rational points of the unit circle. In particular no non-stretchable type
arises at $n=6$.
\end{corollary}

The value $1814$ is due to Jon Hart, who entered it in \texttt{A272906} on 9 May 2016.
What is supplied here is an exhaustive argument: the census of
Section~\ref{sec:upper} bounds $a(6)$ from above by $1814$. We are unaware of a previous
exhaustive proof.

\section{What is counted}\label{sec:counted}

Two clauses of the definition decide the answer, and \texttt{A272906}'s own first comment
fixes both. It reads, verbatim,
\begin{quote}
``For each such slicing, construct the graph on vertices (pieces of the pizza) connected by
edges (line segments separating two pieces). $a(n)$ gives the number of such graphs up to
isomorphism.''
\end{quote}
First, adjacency is \emph{sharing a segment} --- ``line segments separating two pieces'' ---
not sharing a point; the entry's example confirms the strictness by recording that at $n=3$
the five slicings ``exclude a sixth configuration found when the three chords meet at a point
strictly internal to the pizza''. Second, what is counted is \emph{graphs} up to
isomorphism and not arrangements: two arrangements with isomorphic region graphs are one
object. The second clause is why two totals appear below. The census of
Section~\ref{sec:upper} produces $1824$ \emph{arrangement types}; their region graphs fall
into $1814$ \emph{isomorphism classes}, and it is that clause which makes $1814$ rather
than $1824$ the upper bound for $a(6)$. The two numbers first differ at $n=6$: for
$n\le5$ every arrangement type has its own region graph.

\begin{remark}\label{rem:reading}
The prose of \cite{CKS} introduces \texttt{A272906} inside a clause that omits ``in
general position''. Read literally that would permit three chords to be concurrent inside
the disk and would give $6$, not $5$, at $n=3$. \texttt{A272906} does not permit it, and
Theorem~\ref{thm:main} is about \texttt{A272906} as \texttt{A272906} defines it. The
degeneracy-permitting variant is strictly larger and is not settled here.
\end{remark}

Two chords may share a boundary endpoint: they then meet on the circle, not strictly
inside, so the configuration is admissible. The census of Section~\ref{sec:upper}
enumerates only arrangements with $2n$ distinct boundary endpoints, and the following
lemma says nothing is lost.

\begin{lemma}\label{lem:perturb}
Let $A$ be a slicing by $n$ chords in which some boundary points carry more than one
endpoint. Then there is a slicing $A'$ by $n$ chords with $2n$ distinct boundary
endpoints, the same number of interior crossings, and an isomorphic region graph.
\end{lemma}

\begin{proof}
It suffices to separate the endpoints at one boundary point $p$ and iterate. Let
$c_1,\ldots,c_j$ be the chords with an endpoint at $p$, with other endpoints
$q_1,\ldots,q_j$ listed in counterclockwise order starting from $p$. Choose an open arc
$\gamma$ immediately counterclockwise of $p$ that contains no endpoint of any chord, and
pick points $p_j,p_{j-1},\ldots,p_1$ in $\gamma$ in that counterclockwise order; replace
$c_i$ by the chord $p_iq_i$. Two chords $c_a,c_b$ with $a<b$ do not cross after the
change: in the new counterclockwise order
$p_j,\ldots,p_1,q_1,\ldots,q_j$, the arc from $p_a$ to $q_a$ contains
$p_{a-1},\ldots,p_1,q_1,\ldots,q_{a-1}$, and neither $p_b$ nor $q_b$ is among them. A
chord not through $p$ has both endpoints outside $\gamma$, so its interleaving with each
$c_i$ is unchanged, hence so is each crossing. Taking $\gamma$ short enough, the crossings
along each $c_i$ keep their order, so every chord keeps its number of sub-segments and
every crossing keeps its four incident sub-segments.

For a slicing with $n$ chords, $X$ interior crossings and $\beta$ distinct boundary
endpoints, Euler's formula gives $V=\beta+X$, $E=\beta+(n+2X)$ and hence
$F-1=1+n+X$ pieces, independently of $\beta$. So $A$ and $A'$ have the same number of
pieces. Each of the $j+1$ sectors at $p$ in $A$ opens into exactly one piece of $A'$, and
the correspondence extends to a bijection of pieces; since no sub-segment is created or
destroyed, two pieces of $A'$ are separated by a segment if and only if the corresponding
pieces of $A$ are. Hence the region graphs are isomorphic.
\end{proof}

\section{The upper bound}\label{sec:upper}

\subsection*{Types and maps}
Encode an arrangement of $n$ chords with $2n$ distinct boundary endpoints
combinatorially. Number the boundary endpoints $0,\ldots,2n-1$ counterclockwise, give the
chords the labels $0,\ldots,n-1$ in order of the first appearance of an endpoint, and set
\begin{align*}
 \mathrm{bw}[i] &= \text{the label of the chord with an endpoint at position } i,\\
 \mathrm{ord}[c] &= \text{the chords crossing } c, \text{ in order along } c
   \text{ from its first endpoint.}
\end{align*}
The pair $(\mathrm{bw},\mathrm{ord})$ is the \emph{topological type} of the arrangement,
taken up to the dihedral action on the boundary circle and relabelling of chords.

A pair $(\mathrm{bw},\mathrm{ord})$ of a word of length $2n$ over $\{0,\ldots,n-1\}$ and a
list of $n$ sequences of labels is called \emph{consistent} when the following four
conditions hold. Write $X$ for half the total length of the sequences $\mathrm{ord}[c]$.
\begin{enumerate}
\item[(C1)] Every label occurs exactly twice in $\mathrm{bw}$.
\item[(C2)] No label repeats inside one $\mathrm{ord}[c]$, no $\mathrm{ord}[c]$ contains
$c$, and $d\in\mathrm{ord}[c]$ if and only if $c\in\mathrm{ord}[d]$.
\item[(C3)] Whenever $d\in\mathrm{ord}[c]$, exactly one of the two positions of $d$ in
$\mathrm{bw}$ lies strictly between the two positions of $c$; that is, $c$ and $d$
interleave on the boundary circle.
\item[(C4)] The map $M(\mathrm{bw},\mathrm{ord})$ defined next has $V-E+F=2$, has exactly
one face bounded by all $2n$ boundary arcs, and has $1+n+X$ further faces.
\end{enumerate}
Here $M(\mathrm{bw},\mathrm{ord})$ is the following map. Its vertices are the $2n$
boundary positions together with one vertex $x_{cd}$ for each pair with
$d\in\mathrm{ord}[c]$. The \emph{station sequence} of a chord $c$ is its first boundary
position, then the vertices $x_{cd}$ for $d$ running through $\mathrm{ord}[c]$ in that
order, then its second boundary position; the \emph{sub-segments} of $c$ are the
consecutive pairs of that sequence. The edges of $M$ are the $2n$ boundary arcs together
with the $n+2X$ sub-segments. The rotation at a boundary position is (outgoing arc, chord
sub-segment, incoming arc), and the rotation at $x_{cd}$ alternates the two sub-segments of
$c$ with the two of $d$, in the cyclic order fixed by the interleaving of (C3). Faces are
traced from this rotation system in the usual way. A consistent pair determines the pieces
--- the $1+n+X$ faces other than the outer one --- and hence the region graph, in which
two pieces are joined when a sub-segment separates them. The type of a genuine
arrangement is consistent, $M$ being its own map.

\subsection*{The insertion enumeration}
Let $\Tn_1$ consist of the single one-chord type. Let $A$ be a consistent type on $i$
chords. An \emph{insertion datum} for $A$ is a triple
$(\alpha,(\sigma_1,\ldots,\sigma_r),\beta)$ in which
\begin{itemize}
\item $\alpha$ is one of the $2i$ boundary arcs of $A$, and $f_0$ is the piece incident
with $\alpha$;
\item $\sigma_1,\ldots,\sigma_r$ are sub-segments of $A$ and $f_1,\ldots,f_r$ are pieces
with $\sigma_s$ separating $f_{s-1}$ from $f_s$ for $1\le s\le r$;
\item $f_0,\ldots,f_r$ are pairwise distinct, and the chords $c_1,\ldots,c_r$ carrying
$\sigma_1,\ldots,\sigma_r$ are pairwise distinct;
\item $\beta$ is a boundary arc incident with $f_r$.
\end{itemize}
Both $r=0$ and $\beta=\alpha$ are allowed. The datum produces a pair on $i+1$ chords:
insert the new label $i$ into $\mathrm{bw}$ immediately after the arc $\alpha$ and
immediately after the arc $\beta$ --- twice after $\alpha$ if $\beta=\alpha$ ---; for each
$s$ insert the new label into $\mathrm{ord}[c_s]$ at the index of $\sigma_s$ along $c_s$,
so that the new crossing falls where $\sigma_s$ was; and let the sequence of the new chord
be $c_1,\ldots,c_r$, reversed if $\beta$ precedes $\alpha$ on the boundary circle. Then
$\Tn_{i+1}$ is the set of pairs so produced, from every member of $\Tn_i$ and every
insertion datum, \emph{retained only when consistent}, and taken up to the dihedral action
on the boundary circle and relabelling of chords.

\begin{lemma}\label{lem:superset}
The topological type of every $n$-chord slicing with $2n$ distinct boundary endpoints
lies in $\Tn_n$.
\end{lemma}

\begin{proof}
Induction on $n$, the case $n=1$ being $\Tn_1$. Let $A$ be such a slicing with $n\ge2$
chords, let $c$ be one of them and let $A^-=A\setminus\{c\}$, again a slicing with
$2(n-1)$ distinct boundary endpoints, so the type of $A^-$ lies in $\Tn_{n-1}$ by
induction. Two distinct straight lines meet at most once, so $c$ crosses each chord of
$A^-$ at most once. Every piece of $A^-$ is the intersection of the disk with finitely
many half-planes, hence convex, so the line of $c$ meets it in a connected set and $c$
meets each piece of $A^-$ in at most one segment. Therefore the pieces of $A^-$ met by
$c$, listed in order along $c$, are pairwise distinct; the sub-segments of $A^-$ crossed
by $c$, in that same order, lie on pairwise distinct chords; and the two endpoints of $c$
lie in the interiors of two boundary arcs of $A^-$. This is an insertion datum for the
type of $A^-$, and the pair it produces is the type of $A$, which is consistent because
$A$ has a map. Since the enumeration discards only pairs that are not consistent, the type
of $A$ lies in $\Tn_n$.
\end{proof}

Lemma~\ref{lem:superset} makes $\Tn_n$ a superset of the straight-realizable types, so
counting region graphs over $\Tn_n$ can only overshoot $a(n)$. Running the process gives
\[
 |\Tn_2|=2,\quad |\Tn_3|=5,\quad |\Tn_4|=19,\quad |\Tn_5|=130,\quad
 |\Tn_6|=1824,
\]
and, writing $X$ for the number of interior crossings,

\medskip
\noindent
{\small\setlength{\tabcolsep}{2.6pt}
\begin{tabular}{l|*{16}{r}}
$X$ & 0 & 1 & 2 & 3 & 4 & 5 & 6 & 7 & 8 & 9 & 10 & 11 & 12 & 13 & 14 & 15\\
\hline
types & 12 & 30 & 51 & 87 & 108 & 142 & 165 & 179 & 184 & 187 & 182 & 148 &
129 & 101 & 76 & 43\\
graphs & 11 & 28 & 47 & 85 & 107 & 142 & 165 & 179 & 184 & 187 & 182 & 148 &
129 & 101 & 76 & 43
\end{tabular}}

\medskip
\noindent
The second row sums to $1814$. Exactly ten types of $\Tn_6$ have a region graph
isomorphic to that of an earlier type, and all ten collisions occur at $X\le4$. For
$n\le5$ there is no collision at all, and the four totals $2,5,19,130$ agree with the
confirmed terms of \texttt{A272906}. Hence
\begin{equation}\label{eq:upper}
 a(6)\le 1814 .
\end{equation}

\section{The lower bound}\label{sec:lower}

\subsection*{Exact objects} Put
\begin{equation}\label{eq:param}
 P(t)=\left(\frac{1-t^2}{1+t^2},\ \frac{2t}{1+t^2}\right),
\end{equation}
so that $P(t)$ is a rational point of the unit circle for every rational $t$, exactly, and
$t\mapsto P(t)$ is a bijection from $\mathbb Q$ onto the rational points of the circle
other than $(-1,0)$, strictly monotone counterclockwise. An \emph{object} is a strictly
increasing $12$-tuple $t=(t_0,\ldots,t_{11})$ of rationals together with a perfect
matching $M$ of $\{0,\ldots,11\}$; its chords are the segments $P(t_a)P(t_b)$ for
$(a,b)\in M$. Because $P$ is monotone, the index order \emph{is} the counterclockwise
boundary order; equivalently, every triple $i<j<k$ of the points $P(t_i)$ is positively
oriented, which is an exact determinant test.

Two chords $(a,b)$ and $(p,q)$ with four distinct endpoints cross in the open disk if and
only if exactly one of $p,q$ lies strictly between $a$ and $b$; this is an integer test,
and it agrees with the exact segment-intersection test. Writing $A=P(t_a)$, $B=P(t_b)$,
$C=P(t_p)$, $D=P(t_q)$ and $[u,v]=u_1v_2-u_2v_1$, the crossing sits at the parameter
\begin{equation}\label{eq:s}
 s=\frac{[\,C-A,\ D-C\,]}{[\,B-A,\ D-C\,]}\ \in(0,1)
\end{equation}
along the chord $(a,b)$, an exact rational. On each chord the parameters of its crossings
are pairwise distinct if and only if no two of them coincide, and that \emph{is} the
no-strictly-interior-multiple-point clause: three concurrent chords would put two
crossings of one chord at the same $s$. Sorting the parameters along each chord yields
$\mathrm{ord}$, and $(\mathrm{bw},\mathrm{ord})$ is the topological type of the object.

\begin{example}\label{ex:wit}
Take $M=\{(0,6),(1,7),(2,8),(3,9),(4,10),(5,11)\}$ and
\[
 t=(-7,\,-5,\,-4,\,-3,\,-2,\,-1,\,2,\,3,\,4,\,7,\,8,\,9),
\]
row~$1$ of Table~\ref{tab:one}. Every one of the $\binom62=15$ pairs of indices
interleaves, and by \eqref{eq:s} the five crossing parameters on each chord are pairwise
distinct rationals in $(0,1)$, so the object is a slicing with the maximal $X=15$ interior
crossings and $7+15=22$ pieces.
\end{example}

\subsection*{The sweep} Table~\ref{tab:one} lists $47$ objects' worth of boundary data:
$47$ strictly increasing $12$-tuples of integers, row~$1$ being the tuple of
Example~\ref{ex:wit}. For each row and each of the $11!!=10395$ perfect matchings of
$\{0,\ldots,11\}$, form the object and discard it when two crossings coincide on a chord.
Of the $47\cdot10395=488565$ pairs, $487831$ survive and $734$ are discarded. Every
surviving type lies in $\Tn_6$, as Lemma~\ref{lem:superset} demands, and together they
realize \emph{all} $1824$ types of $\Tn_6$. Consequently all $1814$ region graphs of
$\Tn_6$ occur:
\begin{equation}\label{eq:lower}
 a(6)\ge 1814 .
\end{equation}

\begin{proof}[Proof of Theorem~\ref{thm:main} and Corollary~\ref{cor:stretch}]
By Lemma~\ref{lem:perturb} every slicing by six chords has the region graph of one with
$12$ distinct boundary endpoints, and by Lemma~\ref{lem:superset} the type of the latter
lies in $\Tn_6$; the region graphs of $\Tn_6$ fall into $1814$ classes, which is
\eqref{eq:upper}. The sweep exhibits a straight rational slicing for each of the $1824$
types of $\Tn_6$, hence for each of the $1814$ classes, which is \eqref{eq:lower} and, at
the level of types rather than classes, Corollary~\ref{cor:stretch}.
\end{proof}

\begin{table}[htbp]
\footnotesize
\setlength{\tabcolsep}{3.2pt}
\begin{tabular}{r|*{12}{r}}
\multicolumn{1}{c|}{} & \multicolumn{12}{c}{$t_0,\ \ldots,\ t_{11}$}\\
\hline
1 & $-7$ & $-5$ & $-4$ & $-3$ & $-2$ & $-1$ & $2$ & $3$ & $4$ & $7$ & $8$ & $9$ \\
2 & $-60$ & $-50$ & $-39$ & $-37$ & $-36$ & $-20$ & $-3$ & $21$ & $31$ & $39$ & $54$ & $60$ \\
3 & $-59$ & $-51$ & $-18$ & $1$ & $5$ & $6$ & $11$ & $14$ & $19$ & $20$ & $43$ & $47$ \\
4 & $-58$ & $-51$ & $-12$ & $0$ & $3$ & $4$ & $23$ & $31$ & $33$ & $51$ & $58$ & $59$ \\
5 & $-55$ & $-47$ & $-40$ & $-30$ & $-16$ & $-4$ & $0$ & $6$ & $50$ & $51$ & $52$ & $57$ \\
6 & $-55$ & $-40$ & $-39$ & $-36$ & $-31$ & $-28$ & $-26$ & $-18$ & $15$ & $24$ & $57$ & $60$ \\
7 & $-53$ & $-47$ & $-39$ & $-37$ & $-34$ & $-22$ & $-20$ & $4$ & $11$ & $16$ & $27$ & $45$ \\
8 & $-53$ & $-45$ & $-33$ & $-28$ & $-25$ & $-5$ & $17$ & $20$ & $44$ & $53$ & $57$ & $60$ \\
9 & $-51$ & $-45$ & $-44$ & $-36$ & $-35$ & $0$ & $5$ & $27$ & $36$ & $39$ & $50$ & $54$ \\
10 & $-49$ & $-44$ & $-11$ & $-8$ & $13$ & $25$ & $26$ & $27$ & $28$ & $35$ & $46$ & $59$ \\
11 & $-49$ & $-29$ & $-27$ & $-24$ & $-23$ & $-9$ & $21$ & $35$ & $46$ & $54$ & $57$ & $60$ \\
12 & $-40$ & $-38$ & $-18$ & $-9$ & $-1$ & $7$ & $9$ & $14$ & $16$ & $17$ & $24$ & $34$ \\
13 & $-40$ & $-37$ & $-34$ & $-28$ & $-27$ & $-21$ & $3$ & $6$ & $28$ & $32$ & $36$ & $38$ \\
14 & $-40$ & $-36$ & $-35$ & $-27$ & $-22$ & $-14$ & $-13$ & $3$ & $4$ & $7$ & $36$ & $39$ \\
15 & $-40$ & $-36$ & $-32$ & $-23$ & $-21$ & $-15$ & $-14$ & $-12$ & $2$ & $4$ & $21$ & $23$ \\
16 & $-40$ & $-34$ & $-32$ & $-30$ & $-2$ & $8$ & $15$ & $17$ & $19$ & $21$ & $28$ & $40$ \\
17 & $-40$ & $-33$ & $-25$ & $-22$ & $-21$ & $-18$ & $1$ & $20$ & $26$ & $31$ & $37$ & $39$ \\
18 & $-40$ & $-30$ & $-17$ & $-15$ & $-13$ & $-9$ & $-5$ & $2$ & $5$ & $8$ & $20$ & $24$ \\
19 & $-40$ & $-25$ & $-22$ & $-2$ & $0$ & $1$ & $2$ & $3$ & $4$ & $8$ & $10$ & $13$ \\
20 & $-39$ & $-35$ & $-26$ & $-24$ & $-20$ & $-10$ & $-6$ & $3$ & $13$ & $27$ & $30$ & $39$ \\
21 & $-39$ & $-34$ & $-31$ & $-24$ & $-23$ & $-15$ & $-8$ & $3$ & $4$ & $11$ & $15$ & $40$ \\
22 & $-38$ & $-37$ & $-35$ & $-30$ & $-27$ & $-23$ & $-11$ & $6$ & $8$ & $17$ & $31$ & $34$ \\
23 & $-38$ & $-32$ & $-3$ & $2$ & $3$ & $9$ & $16$ & $18$ & $25$ & $26$ & $30$ & $39$ \\
24 & $-38$ & $-29$ & $-27$ & $-23$ & $-19$ & $-17$ & $-12$ & $-9$ & $-8$ & $-1$ & $28$ & $39$ \\
25 & $-38$ & $-29$ & $-20$ & $-15$ & $0$ & $3$ & $5$ & $6$ & $11$ & $13$ & $15$ & $25$ \\
26 & $-38$ & $-28$ & $-14$ & $3$ & $7$ & $12$ & $18$ & $20$ & $25$ & $26$ & $29$ & $33$ \\
27 & $-38$ & $-24$ & $-21$ & $-17$ & $13$ & $16$ & $24$ & $25$ & $26$ & $27$ & $28$ & $34$ \\
28 & $-37$ & $-34$ & $-32$ & $-30$ & $-26$ & $-23$ & $-16$ & $0$ & $8$ & $24$ & $25$ & $38$ \\
29 & $-37$ & $-32$ & $-15$ & $-5$ & $1$ & $16$ & $17$ & $24$ & $25$ & $31$ & $37$ & $38$ \\
30 & $-36$ & $-35$ & $-28$ & $-24$ & $-22$ & $-13$ & $2$ & $13$ & $21$ & $22$ & $25$ & $40$ \\
31 & $-34$ & $-33$ & $-31$ & $-28$ & $-21$ & $-13$ & $1$ & $6$ & $10$ & $24$ & $28$ & $34$ \\
32 & $-34$ & $-29$ & $-25$ & $-14$ & $-1$ & $4$ & $5$ & $8$ & $13$ & $15$ & $20$ & $31$ \\
33 & $-33$ & $-32$ & $-28$ & $-16$ & $-14$ & $7$ & $23$ & $28$ & $30$ & $32$ & $33$ & $39$ \\
34 & $-33$ & $-29$ & $-21$ & $-17$ & $-5$ & $1$ & $11$ & $14$ & $23$ & $25$ & $28$ & $33$ \\
35 & $-33$ & $-25$ & $-17$ & $-14$ & $-8$ & $-5$ & $5$ & $8$ & $10$ & $18$ & $25$ & $38$ \\
36 & $-32$ & $-31$ & $-29$ & $-14$ & $-11$ & $-10$ & $-8$ & $-7$ & $18$ & $20$ & $27$ & $28$ \\
37 & $-32$ & $-31$ & $-28$ & $-22$ & $-19$ & $-11$ & $5$ & $9$ & $11$ & $24$ & $25$ & $35$ \\
38 & $-32$ & $-24$ & $-19$ & $-16$ & $-13$ & $-12$ & $-2$ & $10$ & $11$ & $19$ & $20$ & $27$ \\
39 & $-32$ & $-20$ & $-15$ & $-14$ & $-9$ & $-8$ & $-3$ & $-2$ & $7$ & $30$ & $35$ & $36$ \\
40 & $-31$ & $-30$ & $-28$ & $-27$ & $-20$ & $-14$ & $-7$ & $13$ & $14$ & $17$ & $23$ & $39$ \\
41 & $-31$ & $-25$ & $-19$ & $-4$ & $-2$ & $3$ & $13$ & $17$ & $23$ & $25$ & $27$ & $37$ \\
42 & $-30$ & $-29$ & $-26$ & $-15$ & $-13$ & $-2$ & $16$ & $23$ & $27$ & $30$ & $37$ & $38$ \\
43 & $-30$ & $-29$ & $-24$ & $-20$ & $-19$ & $-18$ & $2$ & $10$ & $11$ & $15$ & $19$ & $21$ \\
44 & $-30$ & $-26$ & $-5$ & $-4$ & $6$ & $7$ & $28$ & $31$ & $32$ & $37$ & $38$ & $39$ \\
45 & $-28$ & $-26$ & $-16$ & $-12$ & $-11$ & $-8$ & $18$ & $22$ & $27$ & $29$ & $30$ & $33$ \\
46 & $-27$ & $-25$ & $-23$ & $-22$ & $-17$ & $-3$ & $-1$ & $13$ & $29$ & $31$ & $33$ & $34$ \\
47 & $-27$ & $-13$ & $-11$ & $-9$ & $-1$ & $10$ & $15$ & $17$ & $20$ & $23$ & $24$ & $29$ \\
\end{tabular}
\medskip
\caption{The $47$ boundary parameter vectors. Swept over all $10395$ perfect matchings
each, they realize all $1824$ topological types of $\Tn_6$ by straight chords with
rational endpoints. Row $1$ is Example~\ref{ex:wit}. Nothing here claims $47$ is
least.}\label{tab:one}
\end{table}

\section{Scope}\label{sec:scope}

Nothing here bears on $a(n)$ for $n\ge7$, nor on the other sequences of Open Problem~2 of
\cite{CKS}. The degeneracy-permitting reading of the definition (Remark~\ref{rem:reading})
is a strictly larger count and is not settled. No claim is made that $47$ parameter
vectors are fewest.

\section{Verification}

Lemma~\ref{lem:perturb} and Lemma~\ref{lem:superset} are hand proofs. The rest is
determined by the printed data: the census of Section~\ref{sec:upper} is fixed by the
insertion rule and takes no input, and the sweep of Section~\ref{sec:lower} is fixed by
Table~\ref{tab:one}.

A program accompanying this note re-derives every computational claim above from the
printed data and nothing else, in exact integer and \texttt{Fraction} arithmetic
throughout, using only the Python standard library. It re-runs the census; counts the
region graphs by two independent routes, a refinement-based canonical form and an
invariant-plus-exact-iso\-morphism search, and checks that the two agree; verifies the
object of Example~\ref{ex:wit}, including a sign-vector derivation of its region graph;
sweeps all $488565$ pairs of Table~\ref{tab:one}; and runs four controls, among them the
reproduction of the confirmed terms $2,5,19,130$. All $53$ of its checks pass. It re-runs
the enumeration Lemma~\ref{lem:superset} licenses, not the lemma.

\begin{thebibliography}{9}

\bibitem{A272906}
J.~Hart, \emph{Number of topologically-distinct pizza slicings from $n$ chords in general
position}, Entry \texttt{A272906} of The On-Line Encyclopedia of Integer Sequences,
9 May 2016. \href{https://oeis.org/A272906}{oeis.org/A272906}. The value $a(6)=1814$ and
the ``unconfirmed'' disclaimer are his; the quotation in Section~\ref{sec:counted} is
comment~1 of the entry.

\bibitem{CKS}
D.~Cutler, A.~Karlsson, and N.~J.~A. Sloane, \emph{Cutting a pancake with an exotic
knife}, arXiv:2511.15864v4, 2026.
\href{https://arxiv.org/abs/2511.15864}{arXiv:2511.15864}.

\end{thebibliography}

\end{document}
